\chapter{Integral Equations}

% -------------------------------------------------------
\section{Definition}
% -------------------------------------------------------

\begin{definition}[Integral Equation]
An \textbf{integral equation} is an equation in which the unknown function appears under an integral sign.
\end{definition}

% -------------------------------------------------------
\section{General Form and Parameters}
% -------------------------------------------------------

The standard form of a linear integral equation is
\[
  y(x) = f(x) + \lambda \int_a^b k(x,t)\,y(t)\,dt,
\]
where:
\begin{itemize}
  \item $y(x)$: the \textbf{unknown function} to be determined,
  \item $f(x)$: the \textbf{free term} (known function),
  \item $k(x,t)$: the \textbf{kernel} (known function of two variables),
  \item $\lambda$: a scalar \textbf{parameter},
  \item $[a,b]$: the limits of integration.
\end{itemize}

% -------------------------------------------------------
\section{Classification}
% -------------------------------------------------------

Integral equations are classified along several independent axes.

\subsection{By Limits of Integration}

\begin{definition}[Fredholm Equation]
An integral equation is of \textbf{Fredholm type} if both limits of integration $a$ and $b$ are \emph{constants}.
\end{definition}

\begin{definition}[Volterra Equation]
An integral equation is of \textbf{Volterra type} if at least one limit of integration is a \emph{variable} (typically the upper limit is $x$).
\end{definition}

\subsection{By Position of the Unknown}

\begin{definition}[First Kind]
The unknown function $y$ appears \emph{only inside} the integral:
\[
  f(x) = \int_a^b k(x,t)\,y(t)\,dt.
\]
\end{definition}

\begin{definition}[Second Kind]
The unknown function $y$ appears both \emph{inside and outside} the integral (i.e.\ $y(x)$ explicitly on the left):
\[
  y(x) = f(x) + \lambda\int_a^b k(x,t)\,y(t)\,dt.
\]
\end{definition}

\subsection{By Linearity}

\begin{definition}[Linear Integral Equation]
The equation is \textbf{linear} if $y$ appears to the first power inside the integral.
\end{definition}

\begin{definition}[Non-Linear Integral Equation]
The equation is \textbf{non-linear} if $y$ appears in a power greater than one or inside a non-linear function under the integral.
\end{definition}

\subsection{By the Free Term}

\begin{definition}[Homogeneous / Non-Homogeneous]
If $f(x) \equiv 0$, the equation is \textbf{homogeneous}; otherwise it is \textbf{non-homogeneous}.
\end{definition}

% -------------------------------------------------------
\section{Classification Examples}
% -------------------------------------------------------

\begin{example}
Classify the integral equation
\[
  u(x) = x + \int_0^1 (x-t)^2\, u(t)\,dt.
\]

\begin{solution}
\begin{itemize}
  \item \textbf{Fredholm}: limits $0$ and $1$ are both constants.
  \item \textbf{Second kind}: $u(x)$ appears outside the integral on the left.
  \item \textbf{Linear}: $u(t)$ appears to the first power inside the integral.
  \item \textbf{Non-homogeneous}: free term $f(x) = x \neq 0$.
\end{itemize}
\end{solution}
\end{example}

% -------------------------------------------------------
\section{Verification of a Solution}
% -------------------------------------------------------

\begin{example}
Verify that $u(x) = e^x$ satisfies the Volterra integral equation of the second kind:
\[
  u(x) = 1 + \int_0^x u(t)\,dt.
\]

\begin{solution}
Substitute $u(t) = e^t$ into the right-hand side:
\[
  1 + \int_0^x e^t\,dt = 1 + \left[e^t\right]_0^x = 1 + e^x - e^0 = 1 + e^x - 1 = e^x.
\]
Since this equals the left-hand side $u(x) = e^x$, the solution is verified. $\checkmark$
\end{solution}
\end{example}
